\begin{korollar}
  {Gauß'scher Min--Max--Satz (Gleichheit in der ersten Bed.)}
  {GaussscherMinMaxSatzGleichheitInDerErstenBedingung}
Seien $r, s \in \N$ und 
$X, Y: \haken r \times \haken s \rightarrow L_2(\Omega, \mfa, \Prim, \R)$
mit 
\[\forall\,i\in\haken r\,\forall\,j\in\haken s\,\exists\,\rho_{i,j}\in\R_{>0}:
X(i, j) \text{ ist }N(0, \rho_{i,j}^2)\text{--verteilt,}\]
\[
\forall\,i\in\haken r\,\forall\,j\in\haken s\,
\exists\,\sigma_{i,j}\in\R_{>0}:
Y(i, j) \text{ ist }N(0, \sigma_{i,j}^2)\text{--verteilt.}\]
Es gelte
\begin{enumerate}[(a)]
	\item $\forall\, i \in \haken r \, \forall\, k,l \in \haken s:
	\norm{X(i,k)-X(i,l)}_2 = \norm{Y(i,k)-Y(i,l)}_2$,
	\item $\forall\, i, j \in \haken r\,\forall\, k, l \in \haken s:
i \neq j \folgt \norm{X(i,k)-X(j,l)}_2 \ge \norm{Y(i,k)-Y(j,l)}_2$.
\end{enumerate}
Dann gilt:
\begin{enumerate}[(1)]
	\item\label{GaussscherMinMaxSatzGleichheitInDerErstenBedingung1}
	$\E\left(\underset{{i \in \haken r}}\min\,
	         \underset{{j \in \haken s}}\max \,X_{i,j} \right)
	\le \E\left(\underset{{i \in \haken r}}\min\,
	         \underset{{j \in \haken s}}\max \,Y_{i,j} \right)$.
	\item\label{GaussscherMinMaxSatzGleichheitInDerErstenBedingung2}
	$\E\left(\underset{{i \in \haken r}}\max\,
	         \underset{{j \in \haken s}}\max \,Y_{i,j} \right)
  \le\E\left(\underset{{i \in \haken r}}\max\,
	         \underset{{j \in \haken s}}\max \,X_{i,j} \right)$.
\end{enumerate}
Insbesondere also
\begin{align*}
\E\bigg(\min_{i \in \haken r}\max_{j \in \haken s}X_{i,j} \bigg)
&\le\E\left(\min_{i \in \haken r}\max_{j \in \haken s}Y_{i,j} \right)
\le\E\left(\max_{i \in \haken r}\max_{j \in \haken s}Y_{i,j} \right)\\
&\le\E\left(\max_{i \in \haken r}\max_{j \in \haken s}X_{i,j} \right)\text.
\end{align*}
\end{korollar}
\begin{proof}
(1) Offenbar sind die Voraussetzungen von~\ref{GaussscherMinMaxSatz} erfüllt.\\
Damit ist bereits die erste Aussage gezeigt.\\
(2) Seien zunächst $a \in \haken r$ und $b, c \in \haken s$.\\
Wegen (a) ist dann $\norm{X_{a,b}-X_{a,c}}_2 = \norm{Y_{a,b} -Y_{a,c}}_2$, sodass
also gilt 
\begin{align}
\forall\, i, j \in \haken r\,\forall\, k, l \in \haken s:
\norm{X_{i,k}-X_{j,l}}_2 \ge \norm{Y_{i,k}-Y_{j,l}}_2 \tag{$\ast$}
\end{align}
Wegen $\abs{\haken r \times \haken s}= \abs{\haken{rs}}$ gibt es
$\sigma: \haken {rs} \rightarrow \haken r \times \haken s$ bijektiv.\\
Seien $\rho_1: \haken r \times \haken s \rightarrow \haken r,(x, y) \mapsto x$,
$\rho_2: \haken r \times \haken s \rightarrow \haken s, (x, y) \mapsto y$ und
$\tau_i\coloneqq  \rho_i \circ \sigma$ für alle $i \in \haken 2$.\\
Seien nun $r'\coloneqq 1, s'\coloneqq  rs$ sowie
\begin{align*}
X' & :\haken {r'}\times \haken {s'} \rightarrow L_2(\Omega, \Prim, \mfa, \R),
(t,z) \mapsto Y(\tau_1(z), \tau_2(z))\text,\\
Y' & :\haken {r'}\times \haken {s'} \rightarrow L_2(\Omega, \Prim, \mfa, \R),
(t,z) \mapsto X(\tau_1(z), \tau_2(z))\text.
\end{align*}
Dann sind $X', Y'$ offensichtlich punktweise zentriert und
normalverteilt, weil $X, Y$ es sind.\\
Seien $k, l \in \haken {s'}$. Dann gilt:
\begin{align*}
&\norm{X'(1,k)-X'(1, l)}_2 = \norm{Y(\tau_1(k), \tau_2(k))
- Y(\tau_1(l), \tau_2(l))}_2\\
\overset{(\ast)}\le &\,
\norm{X(\tau_1(k), \tau_2(k))- X(\tau_1(l), \tau_2(l))}_2
= \norm{Y'(1,k)-Y'(1, l)}_2
\end{align*}
Wegen $r' = 1$ gilt weiter trivial
\begin{align*}
\forall i, j \in \haken {r'}\forall m, n \in \haken {s'}:
i \neq j \folgt \norm{X'(i,m)-X'(j,n)}_2 \ge \norm{Y'(i,m)-Y'(j,n)}_2\text.
\end{align*}
Damit sind die Voraussetzungen von~\ref{GaussscherMinMaxSatz} auch für
$X'$ und $Y'$ erfüllt, also gilt:
\begin{align*}
&\E\left(\max_{i \in \haken r}\max_{i \in \haken s} Y(i,j) \right)
= \E\left(\max_{(i, j) \in \haken r \times \haken s} Y(i,j) \right)
= \E\left(\min_{i\in \haken {r'}}\max_{z \in \haken {s'}}
Y(\tau_1(z), \tau_2(z))\right)\\
= &\, \E\left(\min_{i\in \haken {r'}}\max_{z \in \haken {s'}}
X'(i, z)\right)
\overset{\ref{GaussscherMinMaxSatz}}\le
\E\left(\min_{i\in \haken {r'}}\max_{z \in \haken {s'}}
Y'(i, z)\right)\\
=&\,\E\left(\min_{i\in \haken {r'}}\max_{z \in \haken {s'}}
X(\tau_1(z), \tau_2(z)) \right)
=\E\left(\max_{z \in \haken {s'}}X(\tau_1(z), \tau_2(z)) \right)\\
= & \, \E\left(\max_{(i,j) \in \haken r \times \haken s} X(i,j) \right)
= \E\left(\max_{i \in \haken r}\max_{i \in \haken s} X(i,j) \right)\text.
\end{align*}
Somit ist auch (2) gezeigt.
\end{proof}